\begin{abstract}
Language agents can accumulate debugging experience, but storing prior trajectories is not enough: a memory can make an agent faster on related tasks or brittle on adjacent failures. We study a low-shot alternative for representing debugging experience: inducing a natural-language \texttt{SKILL.md} from a small number of trajectories and treating it as a trajectory-induced procedural memory artifact. The artifact is constrained to contain trigger conditions, an ordered debugging procedure, and failure modes, and is evaluated as an induced memory object rather than as a hand-tuned prompt.

We introduce an evaluation protocol for low-shot procedural skill transfer in language-agent debugging. The protocol records task schemas, diagnosis-edit-test trajectories, solve outcomes, process-efficiency metrics, and negative-transfer annotations. In our primary same-model hard-smoke run over six verified PyBugHive \texttt{black} formatter tasks, Auto SKILL.md matches the best solve rate among primary methods at 5/6 while reducing tokens per solved task relative to Reflexion memory, length-matched Reflexion, a generic checklist, and no memory. Common-solved comparisons show fewer recorded diagnosis-edit-test cycles than reflection baselines. The same run also exposes one explicit negative-transfer failure, showing why transfer evaluation must measure misapplication rather than only solve rate.

A supplementary same-model \texttt{gpt-5.5} structural-control rerun compares Auto SKILL.md with a format-shuffled version that preserves skill content while disrupting its organization. Both methods solve 6/6 tasks; Auto SKILL.md uses fewer recorded cycles and failed patches, while token efficiency is mixed. These results support a cautious conclusion: trajectory-induced procedural memory can produce measurable process-efficiency differences in low-shot within-family debugging, but the current evidence does not establish universal superiority over reflection memory or prove that the canonical SKILL.md structure is necessary for solvability.
\end{abstract}
